\documentclass[11pt,a4paper]{article}

\usepackage[T1]{fontenc}
\usepackage[utf8]{inputenc}
\usepackage{lmodern}
\usepackage[margin=2.5cm]{geometry}
\usepackage{amsmath,amssymb}
\usepackage{array,tabularx}
\usepackage{enumitem}
\usepackage[hidelinks]{hyperref}

\setlist[itemize]{leftmargin=*,itemsep=3pt,topsep=4pt}
\renewcommand{\arraystretch}{1.2}
\newcolumntype{Y}{>{\raggedright\arraybackslash}X}

\title{PV Fault Forecasting Architecture}
\author{}
\date{}

\begin{document}
\maketitle

\section{Problem Definition}

\subsection*{PV Fault Forecasting}
Given a multivariate time series of SCADA signals from an inverter (DC/AC power, voltage, current, solar irradiance, temperature) up to time $t$, the aim is to predict whether a fault of a specific type (open circuit, short circuit, inverter clipping, overheating, ground fault, etc.) will occur within a future horizon $h$ (e.g., the next 1 h, 6 h, or 24 h).

\subsection*{PV Degradation Analysis}
Using a long-horizon time series (weeks to months) of the same SCADA signals, the rate and pattern of gradual performance degradation of the inverter/array must be estimated, independently of weather effects. Possible outputs are the degradation rate (\% per year), a degradation curve, and an estimate of Remaining Useful Life (RUL).

\subsection*{Relationship Between the Two Tasks}
PV-FF is an early warning system: it detects and predicts discrete fault events in the short term. PV-DA is a health monitoring system: it tracks continuous performance degradation over the long term.

\section{Architecture}

\subsection*{Stage 1 -- CS-LSTM Regression}
\emph{Expected behavior under (context-aware) normal conditions.}

\begin{itemize}
    \item \textbf{Purpose:} Predict the expected AC/DC power under normal conditions. The residual, measured $-$ expected, becomes the anomaly signal for all subsequent stages.
    \item \textbf{Rationale:} [A Comprehensive Framework for Predictive Maintenance in Solar Photovoltaic Systems Using Contextual Deep Model] shows that CS-LSTM, by incorporating contextual features (season, time of day) through contextual weighting, achieves $R^2 = 0.9741$ and improves RMSE by 58.96\% over the baseline LSTM, outperforming ANN, RFR, LSTM, and CNN-LSTM with and without contextual features.
    \item \textbf{Contextual features:} Season, time of day (sin/cos encoding), and seasonal decomposition (STL) residuals [Trend-Based Predictive Maintenance and Fault Detection Analytics for Photovoltaic Power Plants].
\end{itemize}

\subsection*{Stage 2 -- XGBoost: Fault Detection \& Classification}
\begin{itemize}
    \item \textbf{Purpose:} (a) Binary Normal/Fault classification; (b) multiclass classification of fault type; (c) prediction of a \emph{Fault Impending} class (3 timesteps before each fault event, as proposed by Baruah et al.), extended to 1-day-ahead prediction.
    \item \textbf{Rationale:} Baruah et al. (2025) compare DT, RF, kNN, AdaBoost, and MLP; RF achieves the best F1 score of 0.992. Many other papers show that XGBoost is among the top methods for both regression and classification tasks because of faster training, built-in regularization, and native feature importance supporting SHAP, compared with, for example, Random Forests.
\end{itemize}

\subsubsection*{XAI Layer -- SHAP, LIME, DiCE}
\begin{itemize}
    \item Integrated on top of the XGBoost classification stage.
    \item \textbf{SHAP values:} Feature importance per prediction (global and local).
    \item \textbf{LIME:} Local interpretability for individual fault events.
    \item \textbf{DiCE:} Generates multiple counterfactuals for a prediction.
    \item \textbf{Rationale:} XAI-enabled fault detection and diagnosis (FDD) increases operator trust and enables fault explanations that can be used for predictive maintenance tasks.
    \item Critical for real-world deployment.
\end{itemize}

\subsection*{Stage 3 -- HMM: Hidden System States (PV-FF Head)}
\begin{itemize}
    \item \textbf{Purpose:} Model the sequence of states $\{\text{Normal} \rightarrow \text{Pre-fault} \rightarrow \text{Fault}\}$ as a Markov chain. The HMM learns transition probabilities between hidden states from the CS-LSTM residuals.
    \item \textbf{Rationale:} Inverter faults do not appear immediately; they develop in phases. An HMM is suitable for sequential state estimation under uncertainty. The HMM + LSTM combination is powerful for fault prognostics.
\end{itemize}

\subsection*{Stage 3 -- LSTM Forecasting Head (PV-FF Head)}
\begin{itemize}
    \item \textbf{Purpose:} Given the HMM state posteriors and the residuals, predict $P(\text{fault\_type}\mid t+h)$ for the horizons $h \in \{1\text{ h}, 6\text{ h}, 24\text{ h}\}$.
    \item \textbf{Rationale:} Baruah et al. (2025) demonstrate early warning up to 2 hours before a fault.
\end{itemize}

\subsection*{Stage 3 -- STL Decomposition + Reliability Score (PV-DA Head)}
\begin{itemize}
    \item \textbf{Purpose:} Separate the long-term degradation trend from seasonality and weather variability. The Reliability Score
    \[
        RS(t) = \frac{PR_{\mathrm{normalized}}(t)}{PR_{\mathrm{baseline}}}
    \]
    tracks gradual degradation. Remaining Useful Life (RUL) is estimated from the slope of the RS trend.
    \item \textbf{Correlation Drift feature:} $D_t = \lVert R(\omega)_t - R(0)\rVert_F$ measures the deviation of the signals' correlation matrix from the initial baseline. It is an early indication of degradation or a structural fault before power loss occurs, and is included as an additional feature in the DA head.
    \item \textbf{Rationale:} [Stanchev \& Hinov] propose a Health Index framework based on normalized performance metrics for PV degradation tracking. STL (Seasonal-Trend decomposition via LOESS) is an established method for time-series decomposition in energy systems.
\end{itemize}

\section{Data}

\subsection*{NIST Dataset}
\url{https://pvdata.nist.gov/}

\noindent\textbf{Information:}\\
Location: Maryland, USA\\
Time granularity: 1 minute

\medskip
\noindent\begin{tabularx}{\textwidth}{@{}lY@{}}
\hline
\textbf{Location} & \textbf{Description} \\
\hline
Ground & Panels installed on the ground \\
Roof & Panels on the roof of a building \\
Canopy & Panels forming a parking canopy \\
ws\_1, ws\_2 & Weather stations: independent meteorological stations \\
\hline
\end{tabularx}

\begin{itemize}
    \item Canopy data (2015, 2016, 2017, 2018): 2,102,400 rows.
    \item Ground data (2015, 2016, 2017, 2018): 2,102,400 rows.
    \item Roof data (2015, 2016, 2017, 2018): 1,706,400 rows.
\end{itemize}

\subsection*{Weather Data}
\begin{tabularx}{\textwidth}{@{}p{0.49\textwidth}Y@{}}
\hline
\textbf{Column} & \textbf{What it measures} \\
\hline
\texttt{Pyra1\_Wm2\_Avg} & Solar irradiance from 2 pyranometers (W/m$^2$) \\
\texttt{RefCell1\_Wm2\_Avg} & Irradiance from a reference cell \\
\texttt{AmbTemp\_C\_Avg} & Ambient temperature ($^\circ$C) \\
\texttt{SEWSPOAIrrad\_Wm2\_Avg} & Plane-of-array irradiance \\
\texttt{WindSpeedAve\_ms}, \texttt{WindDirAve\_deg} & Wind speed and direction \\
\hline
\end{tabularx}

\subsection*{Inverter Data}
\begin{tabularx}{\textwidth}{@{}p{0.41\textwidth}Y@{}}
\hline
\textbf{Column} & \textbf{What it measures} \\
\hline
\texttt{InvPAC\_kW\_Avg} & Inverter AC output power (kW): target \\
\texttt{InvPDC\_kW\_Avg} & DC input power \\
\texttt{InvVDCin\_Avg} & DC input voltage \\
\texttt{InvIDCin\_Avg} & DC input current \\
\texttt{InvVa/b/c\_Avg} & AC voltages (3 phases) \\
\texttt{InvOpStatus\_Avg} & Inverter operating status (0 = off, other = on) \\
\texttt{InvMainFault\_Max} & Inverter faults: fault label \\
\hline
\end{tabularx}

\subsection*{String Currents (Shunts)}
\begin{tabularx}{\textwidth}{@{}p{0.41\textwidth}Y@{}}
\hline
\textbf{Column} & \textbf{What it measures} \\
\hline
\texttt{ShuntCurrent\_A\_Avg\_1} to \texttt{\_7} & Current from 7 separate panel strings, used to identify a problem in a specific string \\
\hline
\end{tabularx}

\smallskip
{\raggedright\noindent\textbf{Fault flags:} \texttt{InvMainFault\_Max}, \texttt{InvDriveFault\_Max}, \texttt{InvVoltageFault\_Max}, \texttt{InvGridFault\_Max}, \texttt{InvTempFault\_Max}, \texttt{InvSystemFault\_Max}, \texttt{InvSystemWarn\_Max}.\par}

\subsection*{Preprocessing}

\subsubsection*{Operating Mask}
\begin{itemize}
    \item Keep only timesteps with irradiance $\geq 100\ \mathrm{W/m^2}$.
    \item The 100 W/m$^2$ threshold excludes the first and last hour of the day, which are noisy due to the low solar angle and high relative sensor uncertainty.
\end{itemize}

\subsubsection*{Outliers}
Domain-knowledge rules are applied before any statistical filter:
\begin{itemize}
    \item Remove negative values.
    \item Depending on the physical interpretation of the features, specific valid value ranges may apply.
    \item Statistical: IQR-based outlier detection for each feature, separately by season.
\end{itemize}

\subsubsection*{Missing Values}
\begin{itemize}
    \item Features with more than 5\% missing values in the entire dataset: remove the feature.
    \item Features with at most 5\% missing values: forward fill / backward fill.
    \item If a gap exceeds 3 minutes, flag the entire corresponding one-hour window and exclude it from training, but not from testing (to evaluate model robustness under noisy conditions).
\end{itemize}

\subsubsection*{Feature Engineering (Before Labeling)}
\begin{itemize}
    \item Performance Ratio (PR):
    \[
        PR = \frac{\texttt{InvPDC\_kW\_Avg}}{(\texttt{SEWSPOAIrrad\_Wm2\_Avg}/1000)\times P_{\mathrm{nominal}}}.
    \]
    \item Inverter efficiency:
    \[
        \mathrm{InvEff} = \frac{\texttt{InvPAC\_kW\_Avg}}{\texttt{InvPDC\_kW\_Avg}}.
    \]
    \item String imbalance:
    \[
        \frac{\operatorname{std}(\text{ShuntCurrent\_A\_Avg\_1..7})}
             {\operatorname{mean}(\text{ShuntCurrent\_A\_Avg\_1..7})}.
    \]
    An important feature for open-circuit / ground-fault detection.
    \item Correlation Drift: $D_t = \lVert R(\omega)_t - R(0)\rVert_F$ over a rolling 30-minute window.
    \item STL decomposition: trend, seasonality, and residual components as separate features for the DA head.
\end{itemize}

\subsubsection*{Labeling}
The procedure analyzed: LSTM-AE anomaly detection $\rightarrow$ physics/rule-based fault typing $\rightarrow$ PU Learning propagation. Fault flags such as \texttt{InvMainFault\_Max} and \texttt{InvTempFault\_Max} are used to validate the generated labels, not as direct training labels.

\subsubsection*{Prediction Targets}
\begin{itemize}
    \item \textbf{PV-FF:} \texttt{InvPAC\_kW\_Avg}---AC output power is the final deliverable of the inverter and the most useful quantity for grid operators. \texttt{InvPDC\_kW\_Avg} serves as an intermediate target for the CS-LSTM regression stage (baseline modeling).
    \item \textbf{PV-DA:} Rolling 30-day PR trend, with seasonality removed through STL.
\end{itemize}

\subsubsection*{Class Imbalance \& Fault Augmentation}
Faults are rare and are expected to account for less than 5\% of operating time. Techniques to address this include:
\begin{itemize}
    \item \textbf{SMOTE:} Creates synthetic samples in feature space between neighboring fault samples. Applied only to the training set.
    \item \textbf{Fault Injection:} Park et al. (2020).
    \item \textbf{Fault Impending class:} The 3 timesteps before each confirmed fault event receive a separate class, as in Baruah et al. (2025), increasing the number of useful training samples.
\end{itemize}

\subsubsection*{Train / Validation / Test Split}
Proposed chronological split using the available data (2015--2018):
\begin{itemize}
    \item 2015, 2016: Training set.
    \item 2017: Validation set.
    \item 2018: Test set.
\end{itemize}

\noindent\textbf{Rationale for this split:}
\begin{itemize}
    \item The training set includes two full years with complete seasonality (two summer/winter cycles), allowing the CS-LSTM to learn seasonal patterns.
    \item The validation set is one full year (2017), sufficient for hyperparameter tuning with complete seasonal coverage, so that no season is favored.
    \item The test set covers one year, allowing evaluation over the full period. Note: the Roof dataset only extends through April 2018.
\end{itemize}

\subsubsection*{Cross-Site Generalization}
Use either the Roof or the Canopy dataset for cross-site verification.

\section{Training / Validation / Testing}

\subsection*{Architecture Training: Sequential, Stage by Stage}
Training will be carried out sequentially, because each stage depends on the output of the previous one.

\begin{center}
Stage 1 $\rightarrow$ Stage 2 $\rightarrow$ PV-FF Head $\rightarrow$ PV-DA Head
\end{center}

\subsection*{Stage 1 -- CS-LSTM Regression}
\begin{itemize}
    \item Trained first, using only confirmed-normal data (filtered 2015--2016 data).
    \item Objective: minimize RMSE between predicted and actual \texttt{InvPAC\_kW\_Avg}.
    \item After training, the model is frozen and its residuals are used as features for Stage 2.
    \item Optimizer: RMSProp (as in Baruah et al., 2025), stable for time-series LSTM.
    \item Loss: MSE.
\end{itemize}

\subsection*{Stage 2 -- XGBoost Classification}
\begin{itemize}
    \item Trained second, with residuals + engineered features + fault labels from the labeling pipeline as input.
    \item SMOTE is applied before fitting, only to the training fold.
    \item Feature importance through built-in XGBoost functionality, followed by SHAP analysis for the XAI layer.
    \item Hyperparameter tuning: GridSearchCV with 5-fold time-series cross-validation.
\end{itemize}

\subsection*{Stage 3 -- PV-FF Head: HMM + LSTM Forecasting}
\begin{itemize}
    \item The HMM is trained on the XGBoost posteriors to learn state sequences.
    \item LSTM forecasting is trained on the HMM output + original features.
    \item Loss: Binary Cross-Entropy for each fault type (one-vs-rest), weighted for class imbalance.
\end{itemize}

\subsection*{Stage 3 -- PV-DA Head: STL + Reliability Score}
\begin{itemize}
    \item STL decomposition.
    \item Reliability Score trend: linear regression on weather-normalized PR, with confidence intervals.
\end{itemize}

\subsection*{Hyperparameters by Model}

\subsubsection*{CS-LSTM}
\begin{tabularx}{\textwidth}{@{}p{0.53\textwidth}Y@{}}
\hline
\textbf{Hyperparameter} & \textbf{Search range} \\
\hline
Hidden units & 10, 15, 20, 25, 30 \\
Sequence length (window) & 5, 10, 15, 30, 60 min \\
Learning rate & $[10^{-5}, 10^{-3}]$ \\
Dropout & 0.1, 0.2, 0.3 \\
Batch size & 64, 128, 256 \\
Epochs & 50--150 \\
LSTM layers & 1, 2 \\
Polynomial basis degree (context) & 2, 3 \\
Optimizer & RMSProp, Adam \\
\hline
\end{tabularx}

\subsubsection*{XGBoost}
\begin{tabularx}{\textwidth}{@{}p{0.53\textwidth}Y@{}}
\hline
\textbf{Hyperparameter} & \textbf{Search range} \\
\hline
\texttt{n\_estimators} & 100, 200, 300, 500 \\
\texttt{max\_depth} & 4, 6, 8, 10 \\
\texttt{learning\_rate} ($\eta$) & 0.01, 0.05, 0.1, 0.3 \\
\hline
\end{tabularx}

\subsubsection*{HMM (PV-FF Head)}
\begin{tabularx}{\textwidth}{@{}p{0.53\textwidth}Y@{}}
\hline
\textbf{Hyperparameter} & \textbf{Value / Rationale} \\
\hline
Number of hidden states & 3: \{Normal, Pre-fault, Fault\} \\
Emission distribution & Gaussian Mixture (3 components) \\
Convergence threshold & $10^{-4}$ \\
N restarts & 10 (to avoid local minima) \\
\hline
\end{tabularx}

\subsubsection*{LSTM Forecasting Head (PV-FF)}
\begin{tabularx}{\textwidth}{@{}p{0.53\textwidth}Y@{}}
\hline
\textbf{Hyperparameter} & \textbf{Range} \\
\hline
Hidden units & 32, 64, 128 \\
Forecast horizons $h$ & 1 h, 6 h, 24 h (separate output head for each horizon) \\
Dropout & 0.1, 0.2, 0.3 \\
Epochs & 50--100 \\
Loss & Weighted Binary Cross-Entropy \\
\hline
\end{tabularx}

\subsection*{Evaluation Metrics}

\subsubsection*{CS-LSTM -- Regression}
\begin{tabularx}{\textwidth}{@{}p{0.4\textwidth}Y@{}}
\hline
\textbf{Metric} & \textbf{What it measures} \\
\hline
RMSE & Absolute error in kW \\
NRMSE & RMSE / mean($P$) \\
$R^2$ & Percentage of explained variance \\
MAE & Robust to outliers \\
\hline
\end{tabularx}

\subsubsection*{XGBoost -- Classification}
\begin{tabularx}{\textwidth}{@{}p{0.4\textwidth}Y@{}}
\hline
\textbf{Metric} & \textbf{What it measures} \\
\hline
F1-score & Average F1 per class and per fault \\
Precision / Recall & Trade-off between false alarms and missed faults \\
AUC-ROC & Threshold-independent separation \\
Confusion Matrix & Complete picture of misclassifications \\
\hline
\end{tabularx}

\subsubsection*{HMM + LSTM -- Fault Forecasting}
\begin{tabularx}{\textwidth}{@{}p{0.4\textwidth}Y@{}}
\hline
\textbf{Metric} & \textbf{What it measures} \\
\hline
Lead time & How early a fault is predicted (in minutes/hours); target: $\geq 2$ h, as in Baruah et al. \\
F1 & F1-score for each forecasting horizon (1 h, 6 h, 24 h) separately \\
False Alarm Rate & $\mathrm{FP}/(\mathrm{FP}+\mathrm{TN})$; critical for preventing alarm fatigue among operators \\
Early Detection Rate & Percentage of faults detected before they occur \\
\hline
\end{tabularx}

\subsubsection*{PV-DA -- Degradation}
\begin{tabularx}{\textwidth}{@{}p{0.4\textwidth}Y@{}}
\hline
\textbf{Metric} & \textbf{What it measures} \\
\hline
RUL MAE & Mean Absolute Error in the Remaining Useful Life estimate \\
PR trend $R^2$ & How well the linear trend explains the weather-normalized PR \\
\hline
\end{tabularx}

\subsection*{Cross-Validation}
\begin{itemize}
    \item \textbf{CS-LSTM:} Time-series cross-validation with an expanding window. Standard k-fold is not used because it breaks temporal order. For example: Train [2015], Val [2016 Q1]; Train [2015--2016 Q1], Val [2016 Q2]; etc., over 4 folds.
    \item \textbf{XGBoost:} $k=5$, preserving chronological order in each fold. SMOTE is applied within each training fold, not to the validation fold.
    \item \textbf{Early stopping:} Used for both CS-LSTM and the LSTM head.
\end{itemize}

\end{document}
